\section{Introduction}
Repeated elliptic solves arise in parameter studies, inverse problems, and design. A learned solution map can supply an initial field to a conventional numerical solver, whose residual criterion then decides acceptance. This division is attractive: prediction need not achieve the final accuracy, while the numerical correction need not begin from zero. Its success, however, depends on the downstream algorithm, problem spectrum, tolerance, and complete computational cost. A small prediction error is not itself evidence of an efficient solver.

Our starting study, S1, fits a quadratic ridge ensemble to a synthetic one-dimensional diffusion family. At relative residual tolerance $10^{-10}$, both learned and zero initializations use 63 iterations on every one of the 128 development evaluation cases. The mean paired iteration reduction is exactly zero in both the in-distribution and coefficient-shifted splits. This result motivates diagnosis rather than deletion or a retrospective endpoint change. S1's original configuration and raw artifacts are preserved, and its confirmatory configuration is evaluated without retuning.

We formulate a separately versioned successor, S2, around an exact structural property: for a fixed diffusion coefficient and homogeneous boundaries, the solution depends linearly on the forcing. An unconstrained quadratic map of coefficient and forcing coordinates contains terms inconsistent with that property. We therefore fit a source-linear coefficient-polynomial map, test it against an unconstrained regression control, and separate the effect of the starting field from a strong classical preconditioner.

The contribution is a compact, auditable computational study. It supplies explicit source identities, preserved failures, complete case-level evidence, and an executable comparison of accuracy, iteration count, and complete online time. It does not introduce a new family of hybrid solvers, prove a universal learned speedup, or establish superiority on realistic multidimensional PDEs. In one dimension a direct tridiagonal algorithm is an especially important alternative; a learned iterative method must justify itself against that inexpensive baseline.

\section{Related work and novelty boundary}
HINTS combines neural operators and classical relaxation using complementary spectral behavior \cite{hints}. NOWS directly studies neural-operator initial guesses for Krylov methods \cite{nows}. These precedents rule out treating learned initialization alone as our novelty. We do not reproduce their architectures, datasets, or published benchmarks and do not compare their reported timings with ours.

POD-augmented conjugate gradients and subspace recycling already integrate data-derived spaces into linear-system solution \cite{pod}. Deflation theory links convergence to subspace quality and the accuracy of its numerical projections \cite{deflation}. Neural Krylov Iteration further demonstrates that learning can be incorporated into Krylov acceleration \cite{nki}. These methods motivate stronger numerical controls; adding a learned basis would not by itself establish a distinct contribution here.

Recent work separates correctness, learning contribution, and complete numerical utility and demonstrates that initial-error summaries can misrank correction work \cite{transfer}. This is directly relevant prior art, not a claim that our study originates that distinction. Reliability analyses of hybrid neural PDE solvers also emphasize physical residuals rather than merely vanishing learned updates \cite{reliable}. We accordingly certify each returned solution using a freshly computed physical residual. Our source-linear architecture is an elementary structural regression model; its value, if any, must be empirical and limited to the declared family.
